\chapter{実験結果}
\section{モデルA--Dの主比較}
表\ref{tab:maincompare}は共通validationの確定済み結果である。A/Bは第1回、Cは第2回、Dは第3回epoch 5である。絶対指標は有効画素、local14は全有効ペアで集計した。
\begin{table}[H]\centering\scriptsize
\caption{共通validationにおける主要モデル比較}\label{tab:maincompare}
\begin{tabular}{lrrrrrr}\toprule
モデル&epoch&MAE m&RMSE m&AbsRel&Pearson $r$&local14 MAE m\\\midrule
A Relative由来Metric&1&2.4897&3.9063&.12120&.8712&.10545\\
B 未学習Metric Outdoor&0&13.2730&14.5340&.65896&.7848&.06561\\
C Metric Outdoor追加学習&1&1.1324&1.7490&.05208&.9773&.06800\\
D Metric Outdoor独立run&5&.7373&1.0094&.03694&.9919&.06148\\\bottomrule
\end{tabular}\end{table}
Bの局所差MAEがCより僅かに低いことだけで形状復元が優れるとは言えない。第2回ゼロ差分基準0.06421 mに対しBは0.06561 m、Cは0.06800 mである。Dは全ペアlocal14が0.06148 mとなり、この基準を下回った。
\section{第3回epoch別推移}
表\ref{tab:epochs}はDriveの第3回epoch CSVと報告表の値である。GT差20 cm以上の局所MAE、振幅比、符号、画像端・内部は固定標本、local14は全有効ペア指標である。
\begin{table}[H]\centering\scriptsize
\caption{モデルDのepoch別validation結果}\label{tab:epochs}
\begin{tabular}{rrrrrrrrrrr}\toprule
Ep.&SiLog&MAE&RMSE&AbsRel&$r$&local14&GT$\geq$20cm MAE&振幅比&符号&秒\\\midrule
1&.10622&1.0791&1.6299&.05088&.9805&.06537&.35276&.3046&78.28\%&1085.5\\
2&.05288&1.7002&1.9953&.08907&.9884&.06339&.34543&.3263&80.24\%&1044.1\\
3&.03596&1.0492&1.3334&.05726&.9894&.06299&.34708&.3228&80.82\%&1097.8\\
4&.03025&.7744&1.0562&.03913&.9901&.06237&.34495&.3428&80.52\%&1024.4\\
5&.02794&.7373&1.0094&.03694&.9919&.06148&.34090&.3304&81.02\%&1020.9\\\bottomrule
\end{tabular}\end{table}
\noindent\small MAE/RMSE/GT局所MAE/local14はm、振幅比・損失は無次元、Pearsonは$r$である。各epochのVRAM peakはallocated 0.877 GiB、reserved 1.248 GiB。AMP skipはepoch 1が1回、epoch 2--5は0回。\normalsize
\begin{figure}[H]\centering
\begin{tikzpicture}
\begin{groupplot}[group style={group size=2 by 3,vertical sep=.9cm,horizontal sep=1.2cm},
width=.39\textwidth,height=3.8cm,grid=major,grid style={gray!20},
xlabel={Epoch},xmin=1,xmax=5,xtick={1,2,3,4,5},tick label style={font=\scriptsize},label style={font=\small},ylabel style={font=\scriptsize}]
\nextgroupplot[ylabel={RMSE (m)},ymin=0,ymax=2.2]\addplot[blue,mark=*] coordinates{(1,1.6299)(2,1.9953)(3,1.3334)(4,1.0562)(5,1.0094)};
\nextgroupplot[ylabel={AbsRel},ymin=0,ymax=.10]\addplot[teal!80!black,mark=square*] coordinates{(1,.05088)(2,.08907)(3,.05726)(4,.03913)(5,.03694)};
\nextgroupplot[ylabel={$\Delta$-MAE (m)},ymin=.055,ymax=.070]\addplot[orange!85!black,mark=triangle*] coordinates{(1,.06537)(2,.06339)(3,.06299)(4,.06237)(5,.06148)};
\nextgroupplot[ylabel={Amplitude ratio},ymin=.25,ymax=.38]\addplot[red!75!black,mark=diamond*] coordinates{(1,.3046)(2,.3263)(3,.3228)(4,.3428)(5,.3304)};
\nextgroupplot[ylabel={Sign match},ymin=.75,ymax=.83,ytick={.76,.78,.80,.82}]\addplot[purple,mark=*] coordinates{(1,.7828)(2,.8024)(3,.8082)(4,.8052)(5,.8102)};
\end{groupplot}\end{tikzpicture}
\caption{第3回モデルDのepoch曲線。既存epoch CSVの数値から再描画した。RMSE/AbsRelは全validation画素、$\Delta$-MAEは全有効14 pxペア、振幅比・符号一致はGT差20 cm以上の固定標本である。縦軸範囲は指標ごとに異なる。}\label{fig:curves}\end{figure}

epoch 2ではRMSEが1.6299 mから1.9953 m、MAEが1.0791 mから1.7002 m、AbsRelが0.05088から0.08907へ一時悪化した。一方local14は0.06537 mから0.06339 mへ改善し、振幅比・符号一致も上昇した。epoch 3--5で絶対誤差は低下し、epoch 5に最小値となった。複数指標は単調に同時改善していない。
\section{GT差20 cm以上の局所構造}
第2回モデルCでは、GT差20 cm以上の標本平均絶対差分0.44370 m、予測平均振幅約0.1453 m、振幅比0.327であった。該当区間局所MAEは0.35423 m、符号一致77.95\%であり、86.72\%のペアで振幅を過小推定した。第3回Dではepoch 1から5へ振幅比0.3046から0.3304、区間MAE0.35276 mから0.34090 m、符号一致78.28\%から81.02\%へ変化した。改善はあるが単調ではない。CとDの振幅比の差は小さいが、標本抽出条件を厳密に突合できないため対応のある差分比較とは解釈しない。
\begin{figure}[H]\centering\begin{tikzpicture}
\begin{axis}[width=.82\textwidth,height=6cm,ybar,bar width=12pt,
symbolic x coords={small,mid1,mid2,large},xtick=data,xticklabels={0--5 cm,5--10 cm,10--20 cm,$\geq$20 cm},
ylabel={Local delta MAE (m)},ymin=0,ymax=.48,
legend style={at={(.5,1.06)},anchor=south,legend columns=2,font=\scriptsize},grid=major,grid style={gray!20}]
\addplot[fill=gray!60] coordinates{(small,.01617)(mid1,.07304)(mid2,.14112)(large,.44370)};
\addplot[fill=blue!70] coordinates{(small,.07480)(mid1,.09063)(mid2,.12255)(large,.41357)};
\addplot[fill=orange!80] coordinates{(small,.02266)(mid1,.06963)(mid2,.12858)(large,.41684)};
\addplot[fill=green!65!black] coordinates{(small,.03451)(mid1,.07031)(mid2,.10600)(large,.35423)};
\legend{Zero baseline,A Relative由来,B 未学習,C 1 epoch}
\end{axis}\end{tikzpicture}
\caption{第2回のGT差分階級ごとの局所深度差MAE。各値は報告済み固定標本。小差分の多数ペアと20 cm以上のペアでモデル順位が異なる。}\label{fig:bins}\end{figure}

第2回GT差分階級標本の72.96\%は0--5 cmであった。同階級の局所MAEはゼロ基準0.01617 mに対しA 0.07480 m、B 0.02266 m、C 0.03451 mである。20 cm以上ではゼロ基準0.44370 m、A 0.41357 m、B 0.41684 m、C 0.35423 mとなり、Cは大きい差分を小さくする誤差を相対的に減らした。
\section{画像端と画像内部}
\begin{figure}[H]\centering\begin{tikzpicture}
\begin{axis}[width=.84\textwidth,height=6.2cm,ybar,bar width=15pt,
symbolic x coords={edge,transition,interior},xtick=data,
xticklabels={0--128 px,128--512 px,$\geq512$ px},
ylabel={14 px local delta MAE (m)},ymin=0,ymax=.43,
legend style={at={(.5,1.06)},anchor=south,legend columns=2,font=\scriptsize},grid=major,grid style={gray!20}]
\addplot[fill=blue!70] coordinates{(edge,.38929)(transition,.09682)(interior,.05814)};
\addplot[fill=orange!80] coordinates{(edge,.08717)(transition,.07033)(interior,.05945)};
\addplot[fill=green!65!black] coordinates{(edge,.09189)(transition,.07112)(interior,.06220)};
\addplot[fill=gray!60] coordinates{(edge,.07866)(transition,.06898)(interior,.05934)};
\legend{A Relative由来,B 未学習,C 1 epoch,ゼロ差分}
\end{axis}\end{tikzpicture}
\caption{第2回固定標本における画像端距離別14 px局所差MAE。端0--128 px帯はペアの10.90\%だが、Aの誤差が突出する。}\label{fig:edge}\end{figure}

第2回標本の画像端0--128 pxではA局所MAE 0.38929 m、B 0.08717 m、C 0.09189 mであった。この帯のペア割合は10.90\%だが、A−Bの総局所MAE増分の82.52\%を占めた。端から512 px以上の内部ではA 0.05814 m、B 0.05945 m、C 0.06220 mとなり順位が逆転する。D epoch 5は端0--128 pxのsampled local MAE 0.08095 m、内部512 px以上0.05678 mである。GT差20 cm以上のDでは端MAE0.36135 m、振幅比0.4951、符号84.40\%、内部MAE0.35727 m、比0.2028、符号77.67\%である。標本結果であり因果効果ではない。
\section{作物・撮影方向別}
\begin{table}[H]\centering\small
\caption{モデルD epoch 5の作物別結果}\label{tab:crop}
\begin{tabular}{lrrrr}\toprule 作物&MAE m&RMSE m&AbsRel&local14 MAE m\\\midrule
Corn&1.00694&1.30455&.04711&.08305\\
Oat&.89419&1.13390&.04414&.06255\\
Soybean&.60541&.86185&.03249&.04892\\
Wheat&.44293&.58861&.02403&.05142\\\bottomrule\end{tabular}\end{table}
\begin{table}[H]\centering\small
\caption{モデルD epoch 5の撮影方向別結果}\label{tab:view}
\begin{tabular}{lrrrrr}\toprule 方向&MAE m&RMSE m&AbsRel&Pearson $r$&local14 MAE m\\\midrule
Nadir&.57857&.75396&.03704&.79186&.03649\\
Oblique&.89654&1.21291&.03684&.99173&.08657\\\bottomrule\end{tabular}\end{table}
ObliqueはNadirよりMAE、RMSE、局所14 MAEが大きい。Pearsonの方向差は深度分布レンジや画像混合集計の影響を受けるため、単独で難易度を決めない。作物・方向・深度範囲が均衡化された要因実験ではないため、差を作物固有の原因へ帰属させない。
\section{checkpoint選択}
第3回報告書ではvalidation RMSE最良はepoch 5（1.0094 m）、GT差20 cm以上の標本局所MAE最良もepoch 5（0.34090 m）である。ただし振幅比はepoch 4の0.3428が最高、符号一致と全14 px MAEはepoch 5が最良である。「局所形状最良」は単一の報告済み選択基準に限った呼称であり、総合最良を意味しない。独立testは未使用である。
\begin{figure}[H]\centering
\begin{tikzpicture}\begin{groupplot}[group style={group size=2 by 1,horizontal sep=1.4cm},width=.43\textwidth,height=5.2cm,ybar=12pt,grid=major,grid style={gray!20},symbolic x coords={A,B,C,D},xtick=data,ylabel style={font=\scriptsize}]
\nextgroupplot[ylabel={RMSE (m)},ymin=0,ymax=16]\addplot[fill=blue!65] coordinates{(A,3.9063)(B,14.5340)(C,1.7490)(D,1.0094)};
\nextgroupplot[ylabel={$\Delta$-MAE (m)},ymin=0,ymax=.12]\addplot[fill=orange!75] coordinates{(A,.10545)(B,.06561)(C,.06800)(D,.06148)};
\addplot[black,dashed,thick,forget plot] coordinates{(A,.064212)(D,.064212)};
\end{groupplot}\end{tikzpicture}
\caption{モデルA--Dのvalidation比較。左は全有効画素RMSE、右は全有効14 pxペア局所MAE。破線は第2回ゼロ差分基準0.06421 m。}\label{fig:compare}\end{figure}



